\documentclass[12pt,a4paper]{report}
\usepackage[utf8]{inputenc}  %mettre latin1 à la place de utf8 si vous utilisez latin1
\usepackage[T1]{fontenc}
\usepackage{setspace}
\renewcommand{\baselinestretch}{2}
\usepackage[french, english]{babel}   %langue française et anglais. Si vous utilisez d'autres langues, ajoutez-les ici. Quand vous avez du texte en anglais, il faut ajouter la commande \selectlanguage{english}. Après, rappelez-vous de revenir au français avec \selectlanguage{french}
\usepackage[maxlevel=3]{csquotes}
\usepackage[backend=biber,style=verbose-trad1,isbn=false]{biblatex}
\DefineBibliographyStrings{french}{in={dans},inseries={dans}}
\usepackage[cyr]{aeguill}
\usepackage{geometry}
\geometry{verbose,letterpaper,tmargin=27.5mm,bmargin=27.5mm,lmargin=27.5mm,rmargin=27.5mm}
\usepackage{graphicx}
\usepackage{epigraph}
\setlength\epigraphwidth{13cm}
\usepackage[center,up,labelfont=bf]{caption}
\usepackage{float}
\usepackage{url}
\newcommand{\guil}[1]{«~{#1}~»}    %guillemets 
\newcommand{\guill}[1]{``{#1}''}     %guillements dans les guillemets
\bibliography{bibliographie/bibliographiebdd}



\begin{document}



\input{titres/titre.tex}
\input{titres/garde.tex}

\pagenumbering{roman} % numérotation en chiffres romains
\setcounter{page}{1}

\input{titres/resume.tex}
\input{titres/abstract.tex}
\input{titres/remerciements.tex}

\begin{singlespace}
\selectlanguage{french}
\tableofcontents % table des matières
\listoffigures % si vous avez des images... si vous n'en avez pas, effacez cette ligne
\end{singlespace}

% fin numérotation en chiffres romains
% début numérotation en chiffres arabs3
\pagenumbering{arabic}
\setcounter{page}{1}


%dans les fichier vous trouverez des exemples d'usage des différentes commandes de LaTeX

\input{textes/intro.tex} %exemple d'intro

\input{textes/chap1.tex} %exemple de chapitre. Vous pouvez écrire le texte en un seul fichier ou le diviser en plusieurs - plus facile à gérer. Créez des nouveaux fichiers  dans le dossier textes et ensuite ajoutez ici les lignes \input


\printbibheading %exemple de bibliographie divisée en sections. Pour ajouter des oeuvres non citées,utiliser \nocite

\printbibliography[keyword=pratique,heading=subbibliography,title={Théories littéraires dans les jeux vidéo}]
\printbibliography[keyword=litteraire,heading=subbibliography,title={Narratologie et structuralisme}]

\printbibliography[keyword=jeu,heading=subbibliography,title={\emph{Games studies}}]

% fin numérotation en chiffres arabs


\end{document} 